% Figure: step-index optical fiber. A ray enters within the acceptance cone,
% refracts at the input face, and zig-zags down the core by total internal
% reflection at the core-cladding boundary.
\begin{figure}[htbp]
\centering
\begin{tikzpicture}[>=Stealth, scale=1.0]
% cladding (outer) and core (inner)
\fill[engblue!8] (0,-1.3) rectangle (9,1.3);
\fill[englime!12] (0,-0.7) rectangle (9,0.7);
\draw[thick] (0,-1.3) -- (9,-1.3);
\draw[thick] (0,1.3) -- (9,1.3);
\draw[enggray, thick] (0,-0.7) -- (9,-0.7);
\draw[enggray, thick] (0,0.7) -- (9,0.7);
% input face
\draw[engblack, very thick] (0,-1.3) -- (0,1.3);
% labels
\node[font=\small, anchor=west] at (9.05,0) {core $n_{\text{core}}$};
\node[font=\small, anchor=west] at (9.05,1.0) {cladding $n_{\text{clad}}$};
% acceptance-cone incident ray entering at center-left
\draw[engred, very thick, ->] (-1.9,0.95) -- (-0.05,0.0);
\node[engred, font=\scriptsize] at (-1.5,0.55) {input ray};
% acceptance cone dashed bounds
\draw[engred, dashed] (-1.9,-0.95) -- (0,0);
\draw[engred] (-1.0,-0.5) arc [start angle=-27, end angle=27, radius=1.12];
\node[engred, font=\scriptsize] at (-1.35,0) {$\theta_a$};
% zig-zag inside the core: bounce points along top/bottom of core (y=+-0.7)
% start at (0,0), go up to (1.4,0.7), down to (2.8,-0.7), up to (4.2,0.7)...
\draw[engorange, very thick, ->]
  (0,0) -- (1.4,0.7) -- (2.8,-0.7) -- (4.2,0.7) -- (5.6,-0.7)
       -- (7.0,0.7) -- (8.4,-0.7) -- (9,-0.4);
% normal and TIR angle at one bounce (at (2.8,-0.7))
\draw[dashed, enggray] (2.8,-0.7) -- (2.8,-1.5);
\draw[dashed, enggray] (2.8,-0.7) -- (2.8,0.1);
\node[engorange, font=\scriptsize, anchor=north] at (2.8,-0.72) {$\theta_c$};
\end{tikzpicture}
\caption[Total internal reflection in a step-index fibre]{A ray entering a
step-index fibre within the acceptance cone of half-angle $\theta_a$
refracts at the input face and then strikes the core-cladding boundary at an
angle exceeding the critical angle $\theta_c$, so it reflects totally and
zig-zags down the core with no loss at each bounce. The numerical aperture
$\text{NA} = \sin\theta_a = \sqrt{n_{\text{core}}^2 - n_{\text{clad}}^2}$
sets the largest input angle the fibre will guide.}
\label{fig:vol04ch11:fiber-tir}
\end{figure}
